\chapter{Introdução}

Geradores de sinais estão entre os instrumentos mais usados em um laboratório de eletrônica: alimentam circuitos sob teste com senoides, rampas e pulsos de frequência e amplitude conhecidas. Os geradores de funções e de formas de onda arbitrárias atuais são, em sua maioria, digitais. Em vez de um oscilador analógico, eles calculam as amostras do sinal e as convertem em tensão por meio de um \ac{dac}.

A técnica que viabiliza esses instrumentos é a \ac{dds}, descrita por \textcite{tierney1971}. Um acumulador de fase soma, a cada ciclo de clock, um valor constante chamado palavra de sintonia; a fase acumulada endereça uma tabela que contém um período da forma de onda; e o \ac{dac} converte as amostras lidas em um sinal analógico, que um filtro passa-baixas suaviza. A frequência de saída é proporcional à palavra de sintonia, o que dá à \ac{dds} resolução de frequência muito fina, troca de frequência sem descontinuidade de fase e formas de onda definidas apenas pelo conteúdo de uma tabela \cites{ad1999,kester2008}.

Os \acp{fpga} são adequados para implementar a \ac{dds}: oferecem memória interna para as tabelas, multiplicadores dedicados, \acp{pll} para gerar o clock e lógica suficiente para acrescentar interfaces de controle. O curso de Engenharia Eletrônica da \ac{utfpr} dispõe da placa de desenvolvimento DE2-115, com um \ac{fpga} Cyclone IV E, que reúne esses recursos \cite{terasic2010}.

Este trabalho apresenta o desenvolvimento de um gerador de formas de onda por \ac{dds} completo, do código em \ac{vhdl} ao conector de saída. O núcleo digital roda no \ac{fpga} da DE2-115 e é controlado pelo computador através do próprio cabo de gravação da placa, usando o recurso Virtual \ac{jtag} da Intel. Uma placa analógica própria converte as amostras em tensão, com um \ac{dac} de 8~bits, um filtro de reconstrução e ajuste de amplitude. Uma interface gráfica em Python permite escolher a frequência, a forma de onda e o conteúdo de uma tabela arbitrária.

Os capítulos seguintes estão organizados da seguinte forma: o Capítulo~\ref{cap:teoria} apresenta a fundamentação teórica da \ac{dds}, dos \acp{fpga} e dos circuitos analógicos de saída; o Capítulo~\ref{cap:metodologia} descreve o projeto do sistema, do núcleo digital, da verificação por simulação, da interface com o usuário e da placa analógica; o Capítulo~\ref{cap:resultados} apresenta e discute os resultados de síntese, de simulação e de bancada; e o Capítulo~\ref{cap:conclusao} reúne as conclusões e as propostas de trabalhos futuros.

\section{Objetivos}

O objetivo principal deste trabalho é desenvolver um gerador de formas de onda por síntese digital direta, implementado em \ac{fpga} e controlado por computador, capaz de gerar senoides, rampas, funções sinc e formas arbitrárias com frequência de 0 a 262\,143~Hz e resolução de 2,33~mHz.

Os objetivos específicos podem ser discriminados como:
\begin{itemize}
	\item Implementar em \ac{vhdl} um núcleo \ac{dds} com acumulador de fase de 32~bits, tabelas de 1024~amostras de 8~bits e conversão da frequência em hertz para a palavra de sintonia sem o uso de divisores;
	\item Implementar o controle pelo computador através do Virtual \ac{jtag}, incluindo a passagem segura dos comandos entre domínios de clock e a escrita da forma arbitrária com o gerador em operação;
	\item Organizar o projeto digital em blocos encapsulados, cada um com interface bem definida, e verificar cada bloco por simulação com \eng{testbenches} auto-verificáveis;
	\item Projetar a placa analógica de saída, com \ac{dac}, conversão corrente-tensão, filtro de reconstrução e ajuste de amplitude, e avaliá-la por simulação SPICE;
	\item Desenvolver uma interface gráfica para configurar a frequência, a forma de onda e a tabela arbitrária;
	\item Validar o sistema completo em bancada.
\end{itemize}

\section{Justificativa}

O presente trabalho se justifica pelos seguintes critérios:
\begin{description}
	\item[Didática] --- em um gerador comercial, o funcionamento interno fica oculto. No sistema desenvolvido, cada etapa da \ac{dds} (acumulação de fase, truncamento, quantização, conversão e reconstrução) pode ser observada, medida e modificada, o que o torna útil para o ensino de processamento digital de sinais e de conversão de dados;
	\item[Flexibilidade] --- como a forma de onda é apenas o conteúdo de uma tabela, o gerador pode reproduzir sinais que um gerador de funções comum não oferece, como o eletrocardiograma sintético usado neste trabalho para testar equipamentos biomédicos;
	\item[Aproveitamento de recursos] --- o sistema usa a placa DE2-115, já disponível no curso, e o controle pelo computador dispensa qualquer hardware além do cabo de gravação. A placa analógica usa componentes comuns, como o \ac{dac} DAC0800 e o amplificador operacional TL074;
	\item[Integração] --- o trabalho reúne projeto digital, projeto analógico, desenvolvimento de software e verificação, competências distribuídas ao longo da formação em Engenharia Eletrônica.
\end{description}
